\documentclass[11pt]{article}

% Camera-ready / submission version, as required by the project spec.
\usepackage[final]{acl}

\usepackage{times}
\usepackage{latexsym}
\usepackage[T1]{fontenc}
\usepackage[utf8]{inputenc}
\usepackage{microtype}
\usepackage{inconsolata}
\usepackage{graphicx}
\usepackage{booktabs}
\usepackage{multirow}
\usepackage{amsmath}
\usepackage{amssymb}
\usepackage{url}

% ---------------------------------------------------------------------------
% Title and authors
% ---------------------------------------------------------------------------
\title{Automated Climate Fact-Checking with Hybrid Retrieval and
NLI-style Evidence Aggregation}

% Replace member names and student IDs before submission.
\author{
  Group 051 \\
  \texttt{Member 1 (Student ID)} \quad
  \texttt{Member 2 (Student ID)} \quad
  \texttt{Member 3 (Student ID)} \\
  School of Computing and Information Systems \\
  The University of Melbourne \\
}

\begin{document}
\maketitle

% ---------------------------------------------------------------------------
\begin{abstract}
We address the COMP90042 climate-claim fact-checking task: for each
claim, retrieve supporting evidence from a corpus of roughly one
million sentences and assign one of four verdicts
(\textsc{Supports}, \textsc{Refutes}, \textsc{Not~Enough~Info},
\textsc{Disputed}). We build a three-stage pipeline. Stage~1 unions the
top-200 results of a fine-tuned BM25 retriever and a fine-tuned BGE
dense retriever; Stage~2 reranks the resulting candidates with a
cross-encoder and keeps the top five evidences; Stage~3 fine-tunes a
DeBERTa-v3-large NLI checkpoint with LoRA in two stages. At inference
we propose \emph{A4.5 hybrid routing}, which trusts a concatenated
4-way classifier on \textsc{Not~Enough~Info} predictions and otherwise
defers to a per-evidence NLI aggregator that can also fire
\textsc{Disputed} when one evidence strongly supports and another
strongly refutes the same claim. On the 154-claim dev set the full
pipeline reaches a claim accuracy of $0.5325$, retrieval F1 of
$0.1057$ and harmonic mean of $0.1764$, beating a concatenated joint
head ($0.5260 / 0.1057 / 0.1761$) and a per-evidence NLI baseline
($0.4675 / 0.1057 / 0.1725$). The headline finding is that routing
buys $+0.65$ pp of accuracy essentially for free; the bottleneck is
retrieval, where the 5-evidence budget caps F1 well below the
77\% Recall@400 achieved by the candidate generator.
\end{abstract}

% ---------------------------------------------------------------------------
\section{Introduction}
% ---------------------------------------------------------------------------

Climate-related misinformation is a high-stakes domain for automatic
fact-checking: claims are often technical, evidence is scattered
across encyclopaedic and scientific text, and a single sentence
rarely settles a question on its own. The shared task gives us
$1{,}208{,}827$ evidence sentences, $1{,}228$ training claims and
$154$ development claims, and asks for both retrieval (a small set of
evidence sentences per claim) and four-way classification of the
overall verdict. We are scored on evidence F1~($F$), claim accuracy
($A$) and their harmonic mean $H = 2FA/(F+A)$.

Three properties make the task hard. First, BM25 alone misses
paraphrastic evidence because climate claims are written in everyday
English while sources use technical wording (e.g., ``melting'' vs.\
``mass-balance loss''). Second, off-the-shelf dense encoders trained
on web search transfer imperfectly to climate prose. Third, the
\textsc{Disputed} class is structurally unrepresented in standard NLI
label spaces, so the strongest off-the-shelf NLI checkpoints cannot
predict it even after fine-tuning.

Our system addresses each of these.
\textbf{(i)} We build a hybrid retriever that takes the union of the
top-200 hits of a fine-tuned BM25 index and a fine-tuned BGE dense
index, lifting Recall@400 from $0.51$ (BM25 alone) to $0.77$.
\textbf{(ii)} We add a cross-encoder reranker that selects the
official 5-evidence subset; this is the only step that has to commit
to a tight evidence budget, and isolates the F1 bottleneck.
\textbf{(iii)} We fine-tune a DeBERTa-v3-large NLI checkpoint with
LoRA in two stages, first on gold evidence and then on retrieved
evidence, so the classifier learns to cope with the noise it will see
at inference.
\textbf{(iv)} As our main original contribution, we propose
\emph{A4.5 hybrid routing}, an inference-time rule that combines two
otherwise complementary heads: a concatenated joint classifier that
is reliable on \textsc{Not~Enough~Info}, and a per-evidence NLI
aggregator that is reliable on \textsc{Supports}/\textsc{Refutes} and
can fire \textsc{Disputed} from a per-evidence support/refute
conflict. The routing rule lifts dev accuracy from $0.5260$ to
$0.5325$ at zero training cost.

% ---------------------------------------------------------------------------
\section{Approach}
% ---------------------------------------------------------------------------

\subsection{Task formulation and notation}
\label{sec:notation}
Let a claim be $c$ and the evidence corpus
$\mathcal{E}=\{e_1,\dots,e_N\}$ with $N\!\approx\!1.2\!\times\!10^6$.
Stage 1 produces a candidate set
$\mathcal{R}_K(c)\!\subseteq\!\mathcal{E}$ of size $K$; Stage 2
reranks $\mathcal{R}_K(c)$ to a fixed 5-element subset
$\mathcal{S}(c)$ that we submit as the evidence prediction. Stage 3
classifies the claim into
$y\!\in\!\{\textsc{Sup},\textsc{Ref},\textsc{Nei},\textsc{Dis}\}$
given $(c,\mathcal{S}(c))$.

\subsection{Stage 1 -- Hybrid retrieval}
\label{sec:retrieval}
\textbf{Sparse.} We use \texttt{BM25Okapi} from \texttt{rank\_bm25}
with the library defaults ($k_1\!=\!1.5$, $b\!=\!0.75$). Documents
and queries are passed through a custom tokeniser that lower-cases,
removes English stop-words, drops punctuation and keeps only
alphanumeric tokens. Evidence shorter than three words or longer than
256 words is filtered out, removing $\approx 4\%$ of the corpus.

\textbf{Dense.} We embed both claims and evidence with
\texttt{BAAI/bge-base-en-v1.5}~\citep{xiao2024bge}
(max length 384, $L_2$-normalised, $d\!=\!768$) and index the
$1.2\!\times\!10^6$ embeddings in FAISS. We then fine-tune the same
encoder with \emph{cached multiple-negatives ranking loss}
on (claim, gold evidence, hard-negative) triplets, where the
hard-negative for each gold pair is mined from the hybrid candidate
set on the training split. The fine-tuned encoder
(\texttt{custom-bge-retriever-final}) replaces the off-the-shelf BGE
in the final pipeline.

\textbf{Fusion.} Rather than a weighted score combination, we use a
\emph{set-level union}: for each claim we retrieve the top-200 hits
from each retriever and take their union, yielding at most 400
candidates,
\begin{equation}
\mathcal{R}_{400}(c) \;=\; \mathrm{Top}_{200}^{\mathrm{BM25}}(c)\;\cup\;
                      \mathrm{Top}_{200}^{\mathrm{BGE}}(c).
\label{eq:hybrid}
\end{equation}
We picked a union rather than a min-max weighted sum because the two
retrievers fail on disjoint claim types (lexical-overlap claims vs.\
paraphrastic ones), and a union avoids the recall loss that any
$\alpha$-blended score would introduce on whichever subset that
retriever already misses. The score-level fusion is deferred to the
reranker. We confirm in \S\ref{sec:results} that this choice closes
the recall gap between BGE and BM25.

\subsection{Stage 2 -- Cross-encoder reranker}
\label{sec:reranker}
We rerank $\mathcal{R}_{400}(c)$ with a cross-encoder built on
\texttt{BAAI/bge-reranker-base}, instantiated as a
\texttt{CrossEncoder} with $\text{num\_labels}\!=\!1$ and
$\text{max\_length}\!=\!384$. Inputs are pairs $(c,e)$ separated by
\verb|[SEP]|. Training pairs are constructed from the train split:
each gold pair $(c,e^{+})$ is paired with four hard negatives mined
by \texttt{sentence\_transformers.util.mine\_hard\_negatives}
($\text{num\_negatives}\!=\!4$, $\text{margin}\!=\!0.1$,
\verb|sampling_strategy="top"|). The model is optimised with binary
cross-entropy for one epoch, batch size~16, learning rate $2{\cdot}10^{-5}$
and warm-up ratio $0.1$. At inference we score every candidate in
$\mathcal{R}_{400}(c)$ and define $\mathcal{S}(c)$ as the five
highest-scored evidences.

\subsection{Stage 3 -- DeBERTa-v3-large + LoRA classifier}
\label{sec:classifier}
\textbf{Backbone.} We start from
\texttt{MoritzLaurer/DeBERTa-v3-large-mnli-fever-anli-ling-wanli}, a
DeBERTa-v3-large~\citep{he2021deberta} checkpoint already adapted to
the entailment family (MNLI, FEVER, ANLI, LingNLI, WANLI). We attach
a fresh 4-way head on top of the \verb|[CLS]| representation. The
backbone is frozen apart from low-rank adapters
\citep{hu2022lora} with $r{=}16$, $\alpha{=}32$, dropout $0.1$,
inserted into the \texttt{query\_proj} and \texttt{value\_proj}
modules; biases are not trained.

\textbf{Two-stage fine-tuning.} Stage~A trains on
\emph{gold} evidence (the evidence sentences linked to each claim in
the training set) for 2 epochs at learning rate $2{\cdot}10^{-4}$.
Stage~B warm-starts from Stage~A and continues for 2 epochs at
$1{\cdot}10^{-4}$ on the top-5 \emph{retrieved} evidence produced by
the reranker, so the classifier is exposed to the same noise it will
face at inference. Both stages use AdamW with cosine warm-up
(warm-up = 10\% of steps), gradient clipping at norm $1.0$, effective
batch size $16$ (per-device $4$, gradient accumulation $4$),
\verb|max_length=512|, and a cross-entropy loss with
inverse-frequency class weights normalised to sum to four, plus
label smoothing $0.1$. The submitted classifier is the Stage~B
checkpoint.

\subsection{Stage 4 -- A4.5 hybrid routing}
\label{sec:routing}
We use two complementary inference modes over the same Stage~3
classifier and combine them at the decision level. The combination
is, to our knowledge, original to this work.

\textbf{A4.1: concatenated joint head.} All five evidences are
concatenated with \verb|[SEP]| separators and passed through the
classifier once; the prediction is the \verb|argmax| over the 4-way
softmax. This path benefits from cross-evidence context and is the
only one that ever fires \textsc{Nei} in practice.

\textbf{A4.3: per-evidence NLI aggregator.} For each
$e_i\!\in\!\mathcal{S}(c)$ we run the classifier on $(c,e_i)$ in
isolation to obtain a 4-way softmax
$p_i=(p_i^{\textsc{sup}}, p_i^{\textsc{ref}}, p_i^{\textsc{nei}}, p_i^{\textsc{dis}})$.
Let $\bar p = \tfrac{1}{|\mathcal{S}(c)|}\sum_i p_i$ and
$\bar p^{(k)} = \max_i p_i^{(k)}$. We assign
\begin{equation}
y^{\text{nli}} \!=\!
\begin{cases}
\textsc{Dis} & \!\!\!\text{if } \bar p^{(\textsc{sup})}\!\geq\!\tau
              \text{ and } \bar p^{(\textsc{ref})}\!\geq\!\tau,\\
\arg\max_k \bar p^{(k)} & \!\!\!\text{otherwise.}
\end{cases}
\label{eq:agg}
\end{equation}
The conflict rule captures the intuition that
\textsc{Disputed} means ``some evidence strongly supports and other
evidence strongly refutes the same claim'', a configuration that no
single softmax can express. The threshold $\tau$ is grid-searched on
dev over $\{0.20, 0.25, \ldots, 0.70, 1.01\}$ ($\tau=1.01$ disables
the rule). In practice, even at the best $\tau$ the aggregator
almost never predicts \textsc{Nei} because the NLI backbone
collapses ``not enough info'' onto neutral evidences.

\textbf{A4.5: routing.} The two heads are then combined by a
two-line rule:
\begin{equation}
y^{\text{final}}\;=\;
\begin{cases}
\textsc{Nei} & \text{if } y^{\text{concat}}=\textsc{Nei},\\
y^{\text{nli}} & \text{otherwise.}
\end{cases}
\label{eq:route}
\end{equation}
That is, we trust the joint head whenever it identifies
\textsc{Not~Enough~Info} (the only class it dominates) and otherwise
let the per-evidence aggregator decide --- including its
\textsc{Disputed} rule. The motivation comes from a per-class
confusion analysis on dev (Table~\ref{tab:perclass}): A4.1 gets
\textsc{Nei} right 27/41 times while A4.3 gets it right 0/41; the
asymmetry is just as stark on \textsc{Supports}/\textsc{Refutes} but
in the opposite direction. Equation~\ref{eq:route} picks the right
head for each claim with no extra parameters.

\subsection{Baselines}
\label{sec:baselines}
We report two within-pipeline ablations that each isolate the
routing rule: \textbf{(B1)} the A4.1 concatenated joint head only,
which is the system without per-evidence aggregation; and
\textbf{(B2)} the A4.3 per-evidence NLI aggregator only, which is
the system without the joint-head fall-back. Both share the same
retrieval and reranker stack as the full system, so $A$ deltas
attribute cleanly to the routing rule of Eq.~\ref{eq:route}.

% ---------------------------------------------------------------------------
\section{Experiments}
% ---------------------------------------------------------------------------

\subsection{Evaluation method}
We evaluate with the official \texttt{eval.py} script, which reports
evidence F1 ($F$), claim accuracy ($A$) and the harmonic mean
$H = 2FA/(F\!+\!A)$ over the 154 dev claims. As a retrieval-side
diagnostic we additionally report Recall@$K$ for
$K\!\in\!\{50, 100, 200, 400\}$, where recall is the fraction of
gold evidences that appear anywhere in the top-$K$ candidate list.
For the classifier we also report a per-class accuracy breakdown,
since the four classes have very different priors ($\approx 44\%$
\textsc{Sup}, $18\%$ \textsc{Ref}, $26\%$ \textsc{Nei}, $12\%$
\textsc{Dis} on dev) and \textsc{Disputed} behaviour is otherwise
hidden by the overall accuracy.

\subsection{Experimental details}
All training runs use Google Colab with a mix of T4 and A100 GPUs.
Heavy artefacts are cached to Google Drive
(\texttt{evidence\_filtered.pkl}, \texttt{bm25.pkl},
\texttt{bge\_evidence\_emb.npy}, \texttt{bge\_faiss.index},
\texttt{hybrid\_\{train,dev,test\}.pkl}) so that retrieval, reranker
training and classifier training can be re-run independently. We
fix a single seed of $42$ throughout. Approximate wall-clock costs
on an A100 are $\approx 25$ min for BGE corpus encoding, $\approx 35$
min for BGE fine-tuning, $\approx 20$ min for reranker training,
and $\approx 60$ min per classifier stage. We did \emph{not} submit
to the test leaderboard; all reported numbers are on dev.

% ---------------------------------------------------------------------------
\section{Results}
\label{sec:results}
% ---------------------------------------------------------------------------

\begin{table}[t]
  \centering
  \small
  \begin{tabular}{lccc}
    \toprule
    \textbf{System} & $F$ & $A$ & $H$ \\
    \midrule
    B1: Joint head only (A4.1)          & 0.106 & 0.526 & 0.176 \\
    B2: NLI aggregator only (A4.3)      & 0.106 & 0.468 & 0.172 \\
    \textbf{Full pipeline (A4.5)}       & \textbf{0.106} & \textbf{0.533} & \textbf{0.176} \\
    \bottomrule
  \end{tabular}
  \caption{Dev-set results. $F$ = evidence F1, $A$ = claim accuracy,
    $H$ = harmonic mean. All three rows share the same retrieval +
    reranker stack and differ only in the classification head, so
    differences in $A$ isolate the effect of the routing rule.}
  \label{tab:main-results}
\end{table}

\begin{table}[t]
  \centering
  \small
  \begin{tabular}{lcccc}
    \toprule
    \textbf{Retriever} & R@50 & R@100 & R@200 & R@400 \\
    \midrule
    BM25                & 0.371 & 0.454 & 0.514 & 0.514 \\
    BGE (off-the-shelf) & 0.507 & 0.617 & 0.683 & 0.683 \\
    Hybrid (Eq.~\ref{eq:hybrid}) & \textbf{0.528} & 0.604 & \textbf{0.713} & \textbf{0.770} \\
    \bottomrule
  \end{tabular}
  \caption{Retrieval Recall@$K$ on dev. BM25 and BGE plateau at $K{=}200$
    because each retriever only emits 200 hits; the hybrid union keeps
    growing because the two retrievers cover disjoint failure modes.}
  \label{tab:recall}
\end{table}

\begin{table}[t]
  \centering
  \small
  \begin{tabular}{lcccc}
    \toprule
    \textbf{Class (n)} & A4.1 & A4.3 & A4.5 & support \\
    \midrule
    \textsc{Supports}~(68)  & 42 & \textbf{57} & 57 & gold \\
    \textsc{Refutes}~(27)   & 12 & \textbf{15} & 15 & gold \\
    \textsc{Nei}~(41)       & \textbf{27} & 0 & 27 & gold \\
    \textsc{Disputed}~(18)  &  0 & 0 &  0 & gold \\
    \bottomrule
  \end{tabular}
  \caption{Per-class correct-count breakdown on the 154-claim dev
    set. A4.1 and A4.3 are strictly complementary on
    \textsc{Sup}/\textsc{Ref} vs.\ \textsc{Nei}; A4.5 inherits the
    best of both. \textsc{Disputed} is missed by all three heads.}
  \label{tab:perclass}
\end{table}

\textbf{Routing pays off.} Table~\ref{tab:main-results} shows that
A4.5 strictly improves over the strongest single head by $+0.65$ pp
of accuracy ($0.5325$ vs.\ $0.5260$) without any extra training.
Table~\ref{tab:perclass} explains why: A4.1 and A4.3 are
\emph{strictly complementary} on \textsc{Nei} vs.\
\textsc{Sup}/\textsc{Ref}, so routing recovers the best decision for
each class.

\textbf{Hybrid retrieval was the right call.}
Table~\ref{tab:recall} shows that the union retriever lifts
Recall@400 from $0.683$ (BGE) and $0.514$ (BM25) to $0.770$. At
$K{=}100$ the two retrievers are roughly tied with BGE marginally
ahead, but the recall gain from union scales as the budget grows,
which is exactly the regime the reranker operates in. The fine-tuned
BGE was the larger lift over off-the-shelf BGE in our exploration
(not shown), so its inclusion in the final system is well motivated
even though we did not have space to ablate it in
Table~\ref{tab:recall}.

\textbf{The bottleneck is retrieval, not classification.} All four
systems in Table~\ref{tab:main-results} share the same
$F\!=\!0.1057$, because the official evidence F1 only credits exact
matches against the 5-evidence budget while gold sets average
$\approx 3.4$ evidences. Even with Recall@400 at $0.77$, the
reranker has to discard 99\% of its candidates, and the gold
evidences are often \emph{semantically} correct paraphrases that
the official F1 does not credit. The harmonic mean is dominated by
this $F$ term, which is why A4.5's accuracy gain translates into a
small $H$ gain. We see this as a retrieval/representation issue
rather than a classifier issue.

\textbf{\textsc{Disputed} remains unsolved.} No system predicts
\textsc{Disputed} correctly on a single dev claim. The conflict rule
in Eq.~\ref{eq:agg} requires \emph{both} a strong-support and a
strong-refute evidence in the same top-5 set, which the reranker
rarely supplies. Lowering $\tau$ trades \textsc{Sup}/\textsc{Ref}
accuracy for sporadic \textsc{Disputed} hits; the grid search picks
the highest~$\tau$ at which dev accuracy still rises, i.e.\
effectively disables the rule. Recovering \textsc{Disputed} would
need a retrieval-side fix (a diversity-aware reranker) or a separate
binary head, neither of which we had time to add.

% ---------------------------------------------------------------------------
\section{Conclusion}
% ---------------------------------------------------------------------------

We built a three-stage retrieval--rerank--NLI pipeline for
climate-claim fact-checking, with one fine-tuned BM25 index, one
fine-tuned BGE encoder, a cross-encoder reranker, a
LoRA-fine-tuned DeBERTa-v3-large classifier, and a novel
inference-time routing rule (A4.5) that combines a joint head and a
per-evidence NLI aggregator at zero extra training cost. On dev the
final system reaches $A=0.5325$ and $H=0.1764$, beating each of the
component heads in isolation and our retrieval baseline.

The main limitations are clear. (i)~Evidence F1 caps the harmonic
mean: with a five-evidence budget and an evaluator that credits only
exact matches, our $F=0.1057$ leaves most of $H$ unrecovered.
(ii)~\textsc{Disputed} is a structural blind spot of NLI-style
classifiers; our conflict-threshold rule cannot fire it when the
reranker rarely puts a support--refute pair in the same top-5.
(iii)~We did not submit to the test leaderboard, so generalisation
to the held-out set is unverified. Promising directions are
diversity-aware reranking to expose conflicting evidence, a
dedicated \textsc{Disputed} binary head trained on synthetic
contradictions, and a soft-evidence F1 that credits paraphrases.

% ---------------------------------------------------------------------------
% The two sections below do NOT count toward the 7-page content limit.
% ---------------------------------------------------------------------------

\section*{Team Contributions}
\begin{itemize}
  \item \textbf{Member 1 (ID)}: hybrid retrieval (BM25 tokeniser,
        BGE encoding, FAISS index, union-fusion evaluation harness),
        Recall@$K$ ablation.
  \item \textbf{Member 2 (ID)}: cross-encoder reranker, hard-negative
        mining, cache infrastructure, end-to-end inference scripts.
  \item \textbf{Member 3 (ID)}: DeBERTa-v3-large + LoRA two-stage
        training, A4.3 NLI aggregator, A4.5 hybrid routing, report
        write-up.
\end{itemize}

% ---------------------------------------------------------------------------
\bibliography{custom}

\end{document}
